\section{Задание 3}

\subsection{Методика измерения}

Для косвенной оценки ассоциативности кеша данных первого уровня (L1d) при помощи программы \texttt{cache\_assoc.cpp} была измерена задержка зависимого обхода строк, адреса которых при обычной индексации относятся к одному набору кеша. Ассоциативность определяет, сколько строк может одновременно находиться в одном наборе. При превышении этого числа строки начинают вытеснять друг друга, поэтому растут доля промахов и задержка обхода.

Память была выровнена по границе страницы; размер страницы получен от ОС и составил 4096 байт. В начале каждой страницы был размещён один 64-байтный узел. По данным \texttt{lscpu}, в L1d имеется 64 набора. При обычной индексации номер набора равен номеру строки по модулю 64, поэтому шаг в $4096 / 64 = 64$ строки сохраняет номер набора. Затрагивалось не более 16 строк (1~КиБ), намного меньше ёмкости L1d. Поэтому скачок задержки связывался с вытеснением внутри одного набора, а не с заполнением всего кеша. Построение и проверка цикла показаны в листинге~\ref{lst:cache-assoc-cycle}.

\verbatimlistingcaption{Размещение узлов, построение и проверка цикла}{lst:cache-assoc-cycle}
\begin{Verbatim}[frame=single,numbers=left,fontsize=\small,breaklines=true]
Алгоритм: построить_цикл_страниц(P, порядок)
Вход: размер страницы P и перестановка индексов узлов.
Выход: память и начальный адрес проверенного цикла.

n <- длина перестановки
память <- n областей по P байт с выравниванием по P байтам
для i от 0 до n - 1:
    узлы[i] <- 64-байтный узел в начале страницы i
для i от 0 до n - 1:
    j <- порядок[i]
    k <- порядок[(i + 1) по модулю n]
    узлы[j].следующий <- адрес(узлы[k])

начало <- адрес(узлы[порядок[0]])
текущий <- начало
для каждого i в порядке перестановки:
    если текущий не равен адрес(узлы[i]): ошибка
    текущий <- текущий.следующий
если текущий не равен началу: ошибка
вернуть память и начало
\end{Verbatim}

Для каждого числа строк было проведено 15 испытаний. Перед измерением выполнялись четыре прогревочных обхода, а измеренное время делилось на число переходов в серии из не менее 5~млн переходов. Обход был выполнен так же, как в листинге~\ref{lst:cache-levels-chase}. После каждого прогревочного обхода и измеряемой серии проверялся возврат к началу. Медианой было восьмое значение упорядоченного списка (листинг~\ref{lst:cache-assoc-measure}).

\verbatimlistingcaption{Измерение задержки обхода одного цикла}{lst:cache-assoc-measure}
\begin{Verbatim}[frame=single,numbers=left,fontsize=\small,breaklines=true]
Алгоритм: измерить_для_числа_строк(начало, n)
Вход: начальный адрес цикла и число узлов n.
Выход: медианная задержка в нс/шаг.

шаги <- n * округлить_вверх(5000000 / n)
задержки <- пустой список
повторить 15 раз:
    повторить 4 раза для прогрева:
        адрес <- обойти(начало, n)
        сохранить адрес и проверить возврат к началу
    t0 <- показание монотонных часов
    адрес <- обойти(начало, шаги)
    t1 <- показание монотонных часов
    сохранить адрес и проверить возврат к началу
    задержка <- (t1 - t0) в наносекундах / шаги
    добавить задержку в задержки
отсортировать задержки по возрастанию
вернуть восьмое значение
\end{Verbatim}

Были проверены циклы из 1--16 строк. Для каждого числа строк перестановка строилась заново, а начальное значение генератора равнялось числу строк (листинг~\ref{lst:cache-assoc-main}).

\verbatimlistingcaption{Перебор числа строк и вывод результатов}{lst:cache-assoc-main}
\begin{Verbatim}[frame=single,numbers=left,fontsize=\small,breaklines=true]
Алгоритм: эксперимент с числом строк в наборе
Вход: границы перебора и размер страницы.
Выход: пары (число строк, медианная задержка в нс/шаг).

P <- размер страницы
если P < 4096, P не кратен 4096 или P не степень двойки:
    завершить работу с ошибкой
для n от 1 до 16:
    инициализировать генератор числом n
    порядок <- перемешать индексы от 0 до n - 1
    память, начало <- построить_цикл_страниц(P, порядок)
    медиана <- измерить_для_числа_строк(начало, n)
    освободить память
    вывести n и медиану
\end{Verbatim}

\subsection{Результаты измерений}

На рисунке~\ref{fig:cache-assoc} представлена зависимость медианной задержки от числа строк, размещённых в одном наборе L1d.

\begin{figure}[htbp]
    \centering
    \begin{tikzpicture}
        \begin{axis}[
                width=0.98\linewidth, height=7cm,
                xlabel={Число строк}, ylabel={Задержка, нс/шаг},
                xtick={1,2,...,16}, grid=major,
                tick label style={font=\small},
                /pgf/number format/use comma,
            ]
            \addplot[blue,thick,mark=*,mark size=1.5pt] coordinates {
                    (1,1.815) (2,1.832) (3,1.864) (4,1.851)
                    (5,1.803) (6,1.818) (7,1.850) (8,1.923)
                    (9,6.194) (10,6.096) (11,6.088) (12,6.442)
                    (13,5.887) (14,5.614) (15,6.097) (16,5.961)
                };
        \end{axis}
    \end{tikzpicture}
    \caption{Медианная задержка при изменении числа строк}
    \label{fig:cache-assoc}
\end{figure}

При восьми строках была получена медианная задержка $1{,}923$~нс/шаг. При добавлении девятой строки она выросла до $6{,}194$~нс/шаг, то есть более чем в три раза. При дальнейшем увеличении числа строк задержка осталась повышенной. Последнее число строк перед скачком было принято за оценку ассоциативности: восемь строк на набор. Для проверки был использован столбец \texttt{WAYS} строки L1d из вывода \texttt{lscpu -C} (таблица~\ref{tab:cache-assoc}).

\begin{table}[H]
    \centering
    \caption{Сравнение оценок ассоциативности L1d}
    \label{tab:cache-assoc}
    \begin{tabular}{|c|c|}
        \hline
        Оценка по графику & Значение по \texttt{lscpu} \\
        \hline
        8 строк на набор  & 8 строк на набор           \\
        \hline
    \end{tabular}
\end{table}
